\begin{proposition}[conductor differential]
\label{prop:conductor-differential}
Let $A$ be a line bundle on $S$ such that $\eta^*A\simeq\cO_{\widetilde S}$,
and choose a nowhere-vanishing trivializing section $e$ of $\eta^*A$.
Put $E=\Omega^1_X|_S\otimes A$.  Assume that $\eta^*\theta_g$ restricts to
zero on $\eta^{-1}(D)$.  Then there is a unique section
\[
 0\ne s\in H^0(S,\Omega^1_X|_S\otimes A)
 \tag{3.4}
\]
whose image in the normalization term of (3.2), namely
$\eta_*(\eta^*E)\oplus E|_D$, is
\[
  \bigl(\eta^*\theta_g\otimes e,\,0\bigr).
 \tag{3.5}
\]
The image of $s$ in $\Omega^1_S\otimes A$ is supported on $D$, hence is a
torsion section and maps to zero in $\widetilde\Omega^1_S\otimes A$.
\end{proposition}

\begin{proof}
Tensor (3.2) by the locally free sheaf
$E=\Omega^1_X|_S\otimes A$.  Finiteness of $\eta$ and the projection formula
identify the first and third terms with $E$ and
$\eta_*(\eta^*E|_{\eta^{-1}(D)})$.  The two components in (3.5) have the
same restriction to $\eta^{-1}(D)$ because the first one is zero there.
Exactness therefore gives a unique section $s$.

On $S\setminus D$, the normalization is an isomorphism and $s$ is
$dg|_S\otimes e$.  Since $g$ is a submersion along the smooth fibre, this is
nonzero on a dense open set, so $s\ne0$.  Its intrinsic image vanishes on
$S\setminus D$ because $g$ is constant on $S$ there.  To see explicitly that
this makes the image torsion, work at a point $x\in S$.  The support of the
germ is contained in the proper closed set $D$; hence, for the finitely
generated module generated by the germ, a power of the ideal of $D$ kills it.
Because $S$ is reduced and no irreducible component of $S$ is contained in
$D$, prime avoidance supplies an element of that ideal outside every minimal
prime of $\mathcal O_{S,x}$.  A suitable power is therefore a non-zero-divisor
annihilating the germ.  The image is torsion and is killed in the quotient by
$\mathscr T_S$.
\end{proof}

